\chapter{Conclusion}

\section{Summary of Contributions}
This work re-implemented the Prune4Web Planner $\to$ Programmatic Element Filter $\to$ Action Grounder pipeline \cite{zhang2025prune4web} on a single 6~GB consumer GPU and extended it along two axes that the original paper does not address. The first extension, DOM Delta Processing, replaces the stateless full-DOM filter with a structural-UID-tracked delta view scored by a hybrid keyword + MiniLM \cite{wang2020minilm} fusion at $\alpha = 0.6$. The second extension introduces a Privacy Hub and a risk-based Human-in-the-Loop (HITL) gate around the grounder, combining a 24-keyword PolicyHub with a regex PII detector and an optional $(\epsilon, 0)$-differential-privacy noise layer for candidate scoring \cite{dettmers2023qlora}.

\section{Empirical Findings}
On all three Mind2Web \cite{deng2023mind2webgeneralistagentweb} test splits (6{,}070 samples in total), DELTA\_HYBRID reduces grounder input by 16 to 18\% in Qwen-tokenised cost while staying within 0.23~pp of Recall@20 on test\_task and test\_website and improving Recall@20 by +1.20~pp on test\_domain, the largest and hardest split (3{,}838 samples). On the end-to-end Element Accuracy metric, a QLoRA-fine-tuned Qwen3-0.6B grounder reaches 88.00\% on Mind2Web test\_task, matching the Prune4Web paper's own fine-tuned Qwen2.5VL-3B-Instruct grounder at 88.28\% within a sub-sample rounding margin while running at roughly 5x fewer parameters and without a vision encoder. Together, these two results show that a sub-1B text-only local grounder, paired with a delta-aware filter, can reproduce the published 3B vision-language baseline at substantially lower compute cost.

\section{Limitations}
Several limitations remain. First, the Element Accuracy numbers in Chapter 5 are measured on 200-sample test\_task and test\_domain slices; the $\pm$6~pp 95\% confidence interval per cell means small differences across grounders should not be over-interpreted. Second, the Privacy Hub and HITL gate have only been exercised qualitatively on live web sessions; no labelled PII benchmark exists in Mind2Web, so quantitative claims about masking precision and HITL trigger rates are deferred to a small pilot. Third, the PrivacyAwareLLM wrapper that injects the \texttt{[PRIVACY DIRECTIVE]} into the system prompt is implemented in \texttt{src/privacy\_aware\_llm.py} and instantiated in \texttt{run\_prune4web.py}, but the planner, filter, and grounder call sites still invoke the raw \texttt{llm\_call} function, so the directive is not actually injected end-to-end in the current driver. Fourth, the DELTA\_HYBRID vs FULL Element Accuracy gap of about 6~pp on the 200-sample slices, although stable across both grounders, points to a downstream sensitivity to candidate ordering that has not yet been controlled for.

\section{Closing Remark}
The combined result of this thesis is that the Prune4Web pipeline can be made cheaper at the filter stage (DOM Delta), safer at the grounder stage (HITL + Privacy Hub), and more accessible at the model stage (sub-1B QLoRA grounder), without sacrificing the headline element-accuracy guarantee of the original paper. The directions in which the system is incomplete, summarised in the next chapter, are concrete and locally tractable rather than open-ended.
